\section{Task Definition}
\label{sec:task}

\paragraph{Setting.}
Let $\mathcal{V}$ be the space of video clips and $\mathcal{L}$ the space of
natural-language strings. CIPHER evaluates a function
$f: \mathcal{V} \times \mathcal{L} \to \Sigma^7$ that maps a corrupted video
$v'$ and an instruction $\mathcal{I}$ to a predicted target string $\hat{k}$.
The original clip $v$ contains a latent key $k^* \in \Sigma^7$ embedded during
a temporal window $[t_s, t_e]$. A sequence of transformation operations
$\mathbf{o} = (o_1, \dots, o_n)$ is applied to $v$ to yield
$v' = o_n \circ \dots \circ o_1 (v)$. The agent's objective is to recover the
inverse program $\mathbf{o}^{-1}$ from $\mathcal{I}$ and $v'$, execute it, and
extract $k^*$ from the recovered terminal state.

The key $k^*$ is drawn from unambiguous alphanumerics
($\Sigma = \text{A--Z} \cup \text{2--9} \setminus \{\text{O,0,I,1}\}$).
Instructions are generated from $\mathbf{o}$ via paraphrase templates.

\paragraph{Three-stage pipeline.}
This task is intentionally composite. An end-to-end failure
($\hat{k} \neq k^*$) may arise because the model inferred the wrong program,
realized the wrong tool chain, or failed to localize and read the key after
successful recovery. By enforcing a three-stage pipeline---Instruction Parsing
(S1), Execution (S2), and Key Extraction (S3)---CIPHER isolates these failure
modes. In policy terms, S1 corresponds to \emph{program inference}, S2 to
\emph{tool-chain realization}, and S3 to \emph{terminal-state extraction}.
This decomposition prevents monolithic metrics from masking which part of the
reasoning policy is actually failing.

\begin{description}
  \item[S1 Parse:] $\mathcal{I} \to \hat{\mathbf{o}}$. Score: sequence-aware op F1.
  \item[S2 Execute:] run \texttt{ffmpeg}($\hat{\mathbf{o}}$). Score: binary validity.
  \item[S3 Extract:] read $\hat{k}$ from edited video. Score: CharF1 + EM.
\end{description}

\paragraph{Operations.}
The operation vocabulary is intentionally minimal: six composable
\texttt{ffmpeg} transforms, \texttt{reverse}, \texttt{hflip},
\texttt{vflip}, \texttt{rotate90}, \texttt{rotate270}, and
\texttt{speed2x}. These are reversible, visually interpretable, and cover both
spatial and temporal recovery. The restricted vocabulary is a design choice:
it keeps the benchmark focused on inverse program recovery rather than on open
ended action-space exploration.

\paragraph{Video-temporal diagnostic domain.}
In addition to the original inverse-\texttt{ffmpeg} pipeline, we instantiate a
controlled video-temporal split-key domain for faster iteration on temporal
aggregation. A key is split into two fragments that appear on two consecutive
frames of a short clip sampled from the public DAVIS-2017 video corpus
\citep{davis2017}. The solver must identify the target frames, read both
fragments, and invert a small symbolic operation before answering. The
operation set is \texttt{none}, \texttt{swap\_fragments},
\texttt{reverse\_chars}, and \texttt{swap\_reverse\_chars}; the last two
require character-level transformation rather than simple OCR.

This domain deliberately separates \emph{temporal evidence acquisition} from
generic video editing: the latent program, target frame indices, fragment
bounding boxes, displayed fragments, normalized fragments, and final answer are
all known from generation. This enables dense process scoring. Besides strict
EM, we assign partial credit for operation recognition, target-frame
localization, displayed-fragment similarity, normalized-answer similarity,
transform-derived answer similarity, final-answer similarity, and successful
tool use. The resulting traces are used both for model diagnosis and for
evolutionary seed selection.
